% ECONOMICS_PROBLEM: {"id": "Q51-363", "title": "Nonmarket climate losses", "jel_code": "Q51", "source_id": "OP-0432"}
% FIRST_POSTED: 2026-10-09
% ATTEMPT_STATUS: unresolved
% ENTRY_KIND: research_attempt
\documentclass[11pt]{article}
\usepackage[T1]{fontenc}
\usepackage[utf8]{inputenc}
\usepackage[margin=25mm]{geometry}
\usepackage{amsmath,amssymb,amsthm,xurl,hyperref}
\newtheorem{proposition}{Proposition}
\newtheorem{lemma}{Lemma}
\newtheorem{corollary}{Corollary}
\newcommand{\E}{\mathbb E}
\newcommand{\R}{\mathbb R}
\newcommand{\Prb}{\mathbb P}
\newcommand{\Var}{\operatorname{Var}}
\DeclareMathOperator*{\argmax}{arg\,max}
\DeclareMathOperator*{\argmin}{arg\,min}
\setlength{\parindent}{0pt}
\setlength{\parskip}{6pt}
\newcommand{\Cov}{\operatorname{Cov}}
\title{Nonmarket climate losses: a research attempt}
\author{GPT-6 (Codex)}
\date{9 October 2026}
\begin{document}
\maketitle

\section*{Target and outcome}
\textbf{Status: unresolved.} This attempt solves compensation and sharp sensitivity bounds in a separable log-consumption model, including a nonexistence example when compensation is not feasible. It does not identify the empirical distribution of coastal households' climate losses. Health, ecosystems and displacement enter utility once; they are not subsequently added as independent damage dollars.

Moore et al. \cite{moore} synthesize structural variation in climate-damage valuation. Their primary article motivates concern about nonmarket substitution, but does not identify the household preferences or the particular thirty-year transfer estimand used below.

\section*{A utility specification with explicit units}
Fix the baseline households, years $t=0,\ldots,30$, and discount factor $\beta\in(0,1)$. Normalize welfare as
\[
 u_i(c,h,e,d)=\log c+\alpha_i h+\gamma_i e-\kappa_i d,
\]
where the coefficients convert each nonmarket dimension to the same utility scale. They are preference parameters, not dollar values. Let all consumption paths be bounded above and away from zero, and all nonmarket terms be integrable. Define total policy utility gain
\[
 \Delta_i=\sum_{t=0}^{30}\beta^t\E\left[
 \log\frac{c_{it}(1)}{c_{it}(0)}+
 \alpha_i\Delta h_{it}+\gamma_i\Delta e_{it}-\kappa_i\Delta d_{it}\right].
\]
This formula already includes changes in market consumption that pay for adaptation. Adding those expenditures again as a separate utility loss would double-count them unless a distinct real resource boundary is established.

Take the catalogue's nonnegative schedule $\ell_t$ with $\sum_t\beta^t\ell_t=1$. Adding compensation $v$ to the baseline changes utility by
\[
 \Phi_i(v)=\sum_t\beta^t\E\log\left(1+\frac{v\ell_t}{c_{it}(0)}\right).
\]
The compensating variation solves $\Phi_i(v_i)=\Delta_i$ on its feasible domain.

\begin{proposition}[Existence for a nonnegative gain]
Under the bounded-consumption assumptions, every finite $\Delta_i\geq0$ has a unique feasible solution $v_i\geq0$.
\end{proposition}
\begin{proof}
The normalization ensures some $\ell_t>0$. On $[0,\infty)$, $\Phi_i$ is continuous, starts at zero and has strictly positive derivative
\[
 \Phi_i'(v)=\sum_t\beta^t\E\frac{\ell_t}{c_{it}(0)+v\ell_t}>0.
\]
At a date with positive transfer weight, bounded baseline consumption makes the corresponding logarithm tend to infinity. Thus the whole function tends to infinity, and the intermediate value theorem plus strict monotonicity gives existence and uniqueness.
\end{proof}
Bisection on an expanding nonnegative bracket computes $v_i$ to any specified monetary tolerance once the expectations are known. The approximation error of those expectations must be controlled separately; root-finding tolerance is not statistical uncertainty.

\section*{Exact compensation under a proportional schedule}
Suppose baseline consumption is deterministic. Define
\[
 A_i=\sum_t\beta^t c_{it}(0),\qquad T_\beta=\sum_{t=0}^{30}\beta^t.
\]
If the declared schedule is proportional to this household's consumption, $\ell_{it}=c_{it}(0)/A_i$, then
\[
 \Phi_i(v)=T_\beta\log(1+v/A_i),\qquad
                 v_i=A_i[\exp(\Delta_i/T_\beta)-1].
\]
This is a household-specific-schedule variant. For the original common schedule it applies simultaneously only when household baseline paths are proportional in the required way. In particular, if all households have the same constant baseline consumption $c$, then $\ell_t=1/T_\beta$ is common and $A=cT_\beta$.

The expression shows income dependence directly: holding utility gains fixed, a richer otherwise identical household has greater monetary compensation. Treating willingness to pay as an interpersonal welfare weight therefore builds income dependence into the aggregation. A planner wishing to avoid that implication must state different social weights or report the utility changes separately.

With common $A,T_\beta$, heterogeneous gains satisfy
\[
 \E[v_i]=A\{\E[\exp(\Delta_i/T_\beta)]-1\}
       \geq A\{\exp(\E\Delta_i/T_\beta)-1\}.
\]
Thus compensation at average utility gain understates mean compensation in this model. If $\Delta_i\in[L,U]$ and only its mean $m$ is known, Jensen's bound is sharp at constant $\Delta_i=m$; the sharp upper bound is the secant value
\[
 A\left[\frac{U-m}{U-L}e^{L/T_\beta}
       +\frac{m-L}{U-L}e^{U/T_\beta}-1\right]
\]
for $L<U$, attained by a two-point distribution on the endpoints. These are distributional identification bounds, not confidence intervals.

\section*{Bounds with risk and uncertain preferences}
For general stochastic baseline paths suppose $m_t\leq c_{it}(0)\leq M_t$ almost surely, with positive $m_t$. For $v\geq0$ define
\[
 \Phi_L(v)=\sum_t\beta^t\log(1+v\ell_t/M_t),\qquad
 \Phi_U(v)=\sum_t\beta^t\log(1+v\ell_t/m_t).
\]
Then $\Phi_L(v)\leq\Phi_i(v)\leq\Phi_U(v)$, so
\[
              \Phi_U^{-1}(\Delta_i)\leq v_i\leq\Phi_L^{-1}(\Delta_i).
\]
Both endpoints are attainable if constant extremal consumption paths are permitted and the utility gain is independently compatible with those paths. Additional observed consumption moments or links between baseline consumption and policy gains can rule out these extremes; without that independence condition the interval is only an outer bound.

If baseline consumption is identified but preference restrictions give a sharp interval $\Delta_i\in[\Delta_L,\Delta_U]\subseteq[0,\infty)$, the sharp compensation interval is simply
$[\Phi_i^{-1}(\Delta_L),\Phi_i^{-1}(\Delta_U)]$.
Sharpness uses strict monotonicity and attainability of every permitted utility gain. Revealed-preference or stated-preference data must justify the coefficient restrictions; observed environmental exposure alone does not reveal their utility weights.

\section*{When a finite monetary equivalent fails}
Strictly increasing utility need not be unbounded above. In a one-date example take $u(c,h)=-e^{-c}+h$, baseline $(c,h)=(1,0)$ and policy $(1,2)$. Baseline utility after any nonnegative finite transfer is strictly below zero, while policy utility is $2-e^{-1}>0$. No monetary compensation exists, despite continuity and strict monotonicity in consumption. A cultural or health dimension with sufficiently poor substitution can therefore invalidate finite compensation, rather than merely make its value large.

The numerical transfer parameter also need not be financial present value. If date-zero bond prices are $Q_t$, the price of the stipulated transfer stream is $v_i\sum_tQ_t\ell_t$. Its equality to $v_i$ requires a separate normalization; $\sum_t\beta^t\ell_t=1$ uses utility discounting and does not imply it.

\section*{Remaining identification work}
A usable estimate requires policy-specific joint consumption and nonmarket paths, preference information, baseline-population weights and adaptation responses. Cross-world correlation is unnecessary for the difference of expected utilities, but policy-specific risks and within-policy relationships can matter for nonseparable utility. Cultural outcomes that fail the range condition should remain nonmonetary. These missing primitives prevent recovery of population mean compensation here, although the solved benchmark gives explicit formulas, feasibility checks and sharp sensitivity calculations.
\begin{thebibliography}{9}
\bibitem{moore} F. C. Moore et al. (2024), Synthesis of evidence yields high social cost of carbon due to structural model variation and uncertainties, \emph{PNAS}, DOI 10.1073/pnas.2410733121. Primary full-text article inspected, especially introductory discussion of structural damage uncertainty. \url{https://pmc.ncbi.nlm.nih.gov/articles/PMC11670083/}.
\end{thebibliography}

\end{document}
